% !TeX root = ../Book5.tex
\begin{figure}[htbp]
\centering
\begin{tikzpicture}[
  x=0.9cm,y=0.9cm,every node/.style={font=\small},
  every path/.style={line width=.6pt},
  root/.style={circle,draw,fill=white,inner sep=0pt,minimum size=4.4pt}]
  % E_6, E_7, E_8 use the coordinate-model numbering; alpha_4 is the fork.
  \node at (2,2.05) {\(E_6\)};
  \draw (0,0)--(4,0) (2,0)--(2,1);
  \foreach \x/\i in {0/1,1/3,2/4,3/5,4/6} {
    \node[root] at (\x,0) {};
    \node at (\x,-0.45) {\(\alpha_\i\)};
  }
  \node[root] at (2,1) {};
  \node at (2,1.45) {\(\alpha_2\)};

  \node at (9.5,2.05) {\(E_7\)};
  \draw (7,0)--(12,0) (9,0)--(9,1);
  \foreach \x/\i in {7/1,8/3,9/4,10/5,11/6,12/7} {
    \node[root] at (\x,0) {};
    \node at (\x,-0.45) {\(\alpha_\i\)};
  }
  \node[root] at (9,1) {};
  \node at (9,1.45) {\(\alpha_2\)};

  \node at (6,-1.7) {\(E_8\)};
  \draw (3,-3.55)--(9,-3.55) (5,-3.55)--(5,-2.55);
  \foreach \x/\i in {3/1,4/3,5/4,6/5,7/6,8/7,9/8} {
    \node[root] at (\x,-3.55) {};
    \node at (\x,-4) {\(\alpha_\i\)};
  }
  \node[root] at (5,-2.55) {};
  \node at (5,-2.1) {\(\alpha_2\)};

  % F_4: alpha_1 and alpha_2 are long; alpha_3 and alpha_4 are short.
  \node at (1.5,-5.1) {\(F_4\)};
  \draw (0,-5.95)--(1,-5.95) (2,-5.95)--(3,-5.95);
  \draw (1,-5.9)--(2,-5.9) (1,-6)--(2,-6);
  \draw (1.37,-5.79)--(1.57,-5.95)--(1.37,-6.11);
  \foreach \x/\i in {0/1,1/2,2/3,3/4} {
    \node[root] at (\x,-5.95) {};
    \node at (\x,-6.4) {\(\alpha_\i\)};
  }

  % G_2: alpha_1 is short, so the triple-edge arrow points left.
  \node at (9.5,-5.1) {\(G_2\)};
  \draw (9,-5.89)--(10,-5.89) (9,-5.95)--(10,-5.95)
    (9,-6.01)--(10,-6.01);
  \draw (9.63,-5.77)--(9.43,-5.95)--(9.63,-6.13);
  \node[root] at (9,-5.95) {};
  \node[root] at (10,-5.95) {};
  \node at (9,-6.4) {\(\alpha_1\)};
  \node at (10,-6.4) {\(\alpha_2\)};
\end{tikzpicture}
\caption[例外型 Dynkin 图]{例外型 Dynkin 图。\(E_6,E_7,E_8\) 的分叉点均为 \(\alpha_4\)，
三个臂长分别为 \((1,2,2)\)、\((1,2,3)\)、\((1,2,4)\)；臂长计中心之外的顶点数。
\(F_4\) 的末两根短，\(G_2\) 的第一根短。}
\label{b5:fig:exceptional-dynkin-diagrams}
\end{figure}
